\documentclass[book,paper=a4,fontsize=10.5pt,oneside,openany]{jlreq}

% 大学物理の問題集 PDF版 共通スタイル

\usepackage[haranoaji]{luatexja-preset}
\usepackage{amsmath,amssymb,mathtools}
\usepackage{tikz}
\usepackage[most]{tcolorbox}
\usepackage{xcolor}
\usepackage{geometry}
\usepackage{fancyhdr}
\usepackage{enumitem}
\usepackage{hyperref}
\usepackage{bookmark}

\geometry{
  a4paper,
  top=23mm,
  bottom=24mm,
  inner=23mm,
  outer=20mm,
  headheight=28pt,
  headsep=5mm
}

\definecolor{PhysicsNavy}{HTML}{151E3D}
\definecolor{PhysicsBlueGray}{HTML}{E9EDF4}
\definecolor{PhysicsRule}{HTML}{AAB3C3}
\definecolor{PhysicsText}{HTML}{20242C}

\color{PhysicsText}
\setlength{\parindent}{1em}
\setlength{\parskip}{0.25\baselineskip}
\linespread{1.12}

\hypersetup{
  colorlinks=true,
  linkcolor=PhysicsNavy,
  urlcolor=PhysicsNavy,
  citecolor=PhysicsNavy,
  bookmarksnumbered=true
}

% 章・節見出し（jlreq純正の設定を利用）
\ModifyHeading{chapter}{
  font={\Huge\bfseries\color{PhysicsNavy}},
  label_format={第\thechapter 章}
}
\ModifyHeading{section}{
  font={\Large\bfseries\color{PhysicsNavy}}
}

% ヘッダー・フッター
\pagestyle{fancy}
\fancyhf{}
\fancyhead[L]{\small\color{PhysicsNavy}\leftmark}
\fancyhead[R]{\small\color{PhysicsNavy}大学物理の問題集}
\fancyfoot[C]{\small\color{PhysicsNavy}\thepage}
\renewcommand{\headrulewidth}{0.4pt}
\renewcommand{\headrule}{\hbox to\headwidth{\color{PhysicsRule}\leaders\hrule height \headrulewidth\hfill}}
\renewcommand{\chaptermark}[1]{\markboth{第\thechapter 章\quad #1}{}}

% 問題ボックス
\newtcolorbox{problemBox}[2][]{
  enhanced,
  breakable,
  colback=white,
  colframe=PhysicsNavy,
  colbacktitle=PhysicsNavy,
  coltitle=white,
  title={#2},
  fonttitle=\bfseries,
  boxrule=0.7pt,
  arc=1.2mm,
  left=5mm,
  right=5mm,
  top=4mm,
  bottom=4mm,
  before skip=0.5\baselineskip,
  after skip=\baselineskip,
  attach boxed title to top left={xshift=0mm},
  boxed title style={
    sharp corners=south,
    boxrule=0pt,
    left=4mm,
    right=4mm,
    top=1.6mm,
    bottom=1.6mm
  },
  #1
}

\newcommand{\solutionheading}{%
  \section*{解説}%
  \addcontentsline{toc}{section}{解説}%
  \noindent\color{PhysicsNavy}\rule{\linewidth}{0.8pt}%
  \par\vspace{0.9\baselineskip}\color{PhysicsText}%
}

\newcommand{\solutionpart}[1]{%
  \par\vspace{1.1\baselineskip}%
  \noindent{\large\bfseries\color{PhysicsNavy}#1}%
  \par\vspace{0.35\baselineskip}\color{PhysicsText}%
}

\setlist[enumerate]{
  label=(\arabic*),
  leftmargin=3em,
  itemsep=0.75\baselineskip,
  topsep=0.6\baselineskip
}

\setcounter{tocdepth}{1}


\hypersetup{
  pdftitle={大学物理の問題集 力学},
  pdfsubject={大学物理の演習問題と解説},
  pdfcreator={LuaLaTeX}
}

\begin{document}

\hypersetup{pageanchor=false}
\begin{titlepage}
  \thispagestyle{empty}
  \begin{tikzpicture}[remember picture,overlay]
    \fill[PhysicsNavy]
      (current page.north west) rectangle
      ([yshift=-12mm]current page.north east);
    \fill[PhysicsNavy]
      ([yshift=12mm]current page.south west) rectangle
      (current page.south east);
    \draw[PhysicsRule,line width=1pt]
      ([xshift=23mm,yshift=-48mm]current page.north west) --
      ([xshift=-23mm,yshift=-48mm]current page.north east);
  \end{tikzpicture}

  \vspace*{38mm}
  {\Large\color{PhysicsNavy}大学物理の問題集\par}
  \vspace{7mm}
  {\fontsize{34pt}{42pt}\selectfont\bfseries\color{PhysicsNavy}力学\par}
  \vspace{5mm}
  {\large 演習問題と解説\par}

  \vfill
  \begin{flushright}
    {\large 第1章\quad 運動学と座標系\par}
    \vspace{3mm}
    {\small 試作版・第1問収録\par}
  \end{flushright}
  \vspace*{22mm}
\end{titlepage}
\hypersetup{pageanchor=true}

\frontmatter
\tableofcontents
\clearpage

\mainmatter
\pagestyle{fancy}
\chapter{運動学と座標系}
\thispagestyle{fancy}

\addcontentsline{toc}{section}{1.1　等加速度運動の一般解}

\begin{problemBox}{問題 1.1\quad 等加速度運動の一般解\hspace{1.2em}{\small\mdseries 〔基本〕}}
質量 $m$ の無人探査機が，宇宙空間を一直線上に運動している。軌道に沿って $x$ 軸をとり，探査機は $x$ 軸の正方向に一定の推進力 $F_0$ を受けるものとする。

時刻 $t=0$ において，探査機は原点 $x=0$ を速度 $v_0>0$ で通過した。探査機を質点とみなし，以下の問いに答えよ。

\begin{enumerate}
  \item 探査機の位置を $x(t)$，速度を
  $v(t)=\dfrac{dx}{dt}$ とするとき，運動方程式を示せ。

  \item (1) で得られた運動方程式を，速度 $v(t)$ に関する1階微分方程式とみなす。初期条件 $v(0)=v_0$ のもとでこの微分方程式を解き，速度 $v(t)$ を時間の関数として求めよ。

  \item (2) の結果と $v(t)=\dfrac{dx}{dt}$ の関係を用いて，位置 $x(t)$ に関する微分方程式を立てよ。これを初期条件 $x(0)=0$ のもとで解き，位置 $x(t)$ を時間の関数として求めよ。

  \item 探査機が原点から距離 $L>0$ だけ進んだときの速度 $v_L$ を，$m$，$F_0$，$v_0$，$L$ を用いて表せ。
\end{enumerate}
\end{problemBox}

\solutionheading
\begin{solutionparts}

\solutionpart{(1)}
運動方程式より，質量と加速度の積は物体に働く外力に等しくなります。加速度は $\dfrac{d^2x}{dt^2}$ または $\dfrac{dv}{dt}$ と表されます。

探査機に働く外力は $x$ 軸正の向きに一定の力 $F_0$ であるため，求める運動方程式は以下のようになります。
\begin{equation*}
  m\frac{d^2x}{dt^2}=F_0
\end{equation*}
あるいは，速度 $v$ を用いて以下のように表すこともできます。
\begin{equation*}
  m\frac{dv}{dt}=F_0
\end{equation*}

\solutionpart{(2)}
(1) で得られた運動方程式の両辺を $m>0$ で割ると，次のようになります。
\begin{equation*}
  \frac{dv}{dt}=\frac{F_0}{m}
\end{equation*}
この方程式は，両辺を時間 $t$ で積分することにより解くことができます。
\begin{equation*}
  v(t)=\int\frac{F_0}{m}\,dt=\frac{F_0}{m}t+C_1
\end{equation*}
ここで $C_1$ は積分定数です。初期条件 $v(0)=v_0$ を代入すると，
\begin{equation*}
  v(0)=\frac{F_0}{m}\cdot 0+C_1=v_0
\end{equation*}
となり，$C_1=v_0$ が得られます。したがって，速度 $v(t)$ の時間依存性は以下のように求まります。
\begin{equation*}
  \boxed{v(t)=v_0+\frac{F_0}{m}t}
\end{equation*}
加速度は一定値 $a=\dfrac{F_0}{m}$ です。

\solutionpart{(3)}
速度と位置の関係式 $v(t)=\dfrac{dx}{dt}$ に (2) の結果を代入すると，次の1階微分方程式が得られます。
\begin{equation*}
  \frac{dx}{dt}=v_0+\frac{F_0}{m}t
\end{equation*}
この方程式の両辺を時間 $t$ で積分します。
\begin{align*}
  x(t)
  &=\int\left(v_0+\frac{F_0}{m}t\right)dt \\
  &=v_0t+\frac{F_0}{2m}t^2+C_2
\end{align*}
ここで $C_2$ は積分定数です。初期条件 $x(0)=0$ を代入すると，
\begin{equation*}
  x(0)=v_0\cdot 0+\frac{F_0}{2m}\cdot 0^2+C_2=0
\end{equation*}
となり，$C_2=0$ が得られます。したがって，位置 $x(t)$ の時間依存性は以下のようになります。
\begin{equation*}
  \boxed{x(t)=v_0t+\frac{F_0}{2m}t^2}
\end{equation*}
これは一定加速度 $a=\dfrac{F_0}{m}$ の等加速度運動を表します。

\solutionpart{(4)}
探査機が距離 $L$ だけ進んだときの時刻を $t_L$ とします。このとき，$x(t_L)=L$ が成り立ちます。(3) の結果より，次の関係式が得られます。
\begin{equation*}
  v_0t_L+\frac{F_0}{2m}t_L^2=L
\end{equation*}
速度と位置の式から時刻 $t$ を消去します。(2) より，時刻 $t$ における速度 $v(t)$ の式から $t$ を孤立させます。
\begin{equation*}
  t=\frac{m}{F_0}\bigl(v(t)-v_0\bigr)
\end{equation*}
この式を (3) の位置の式に代入します。
\begin{align*}
  x(t)
  &=v_0\left\{\frac{m}{F_0}\bigl(v(t)-v_0\bigr)\right\}
    +\frac{F_0}{2m}
    \left\{\frac{m}{F_0}\bigl(v(t)-v_0\bigr)\right\}^2 \\
  &=\frac{mv_0}{F_0}\bigl(v(t)-v_0\bigr)
    +\frac{m}{2F_0}\bigl(v(t)-v_0\bigr)^2 \\
  &=\frac{m}{2F_0}\bigl(v(t)-v_0\bigr)
    \left\{2v_0+\bigl(v(t)-v_0\bigr)\right\} \\
  &=\frac{m}{2F_0}\bigl(v(t)-v_0\bigr)\bigl(v(t)+v_0\bigr) \\
  &=\frac{m}{2F_0}\bigl(v(t)^2-v_0^2\bigr).
\end{align*}
この式を整理すると，以下の関係式が得られます。
\begin{equation*}
  v(t)^2-v_0^2=\frac{2F_0}{m}x(t)
\end{equation*}
探査機が位置 $x(t_L)=L$ に達したときの速度を $v_L=v(t_L)$ とすると，
\begin{equation*}
  v_L^2=v_0^2+\frac{2F_0L}{m}
\end{equation*}
探査機は正の向きの力を受け続けて加速するため，$v_L>0$ です。したがって，平方根の正の値を採用することで，求める速度 $v_L$ は以下のように得られます。
\begin{equation*}
  \boxed{v_L=\sqrt{v_0^2+\frac{2F_0L}{m}}}
\end{equation*}

\end{solutionparts}


\end{document}
